\documentclass{article}
\usepackage{mathnotes}

\title{Probability Basics}
\description{Foundational concepts including sample spaces, events, probability rules, conditional probability, and independence using set theory operations.}

\source{title=Log probability, type=web, url={https://en.wikipedia.org/wiki/Log_probability}, accessed=2026-07-21}
\source{title=Probability, author=Nathan Ponder, type=book, published=2020, notes={Louisiana State University at Alexandria}}

\begin{document}

\section{Probability Basics}

\subsection{Sample Space}

\begin{definition}[Sample Space]
The \@{set} of all possible outcomes of a statistical experiment is called the \textbf{sample space} and is represented by the symbol $S$. Each outcome in a sample space is called an \textbf{\@{element}} or \textbf{\@{member}} of the sample space or simply a \textbf{sample point.}

A sample space must be non-empty.
\end{definition}

For example, the \@{sample space} representing the outcomes of the rolls of a standard six-sided die is

\[ S = \{1, 2, 3, 4, 5, 6\}. \]

\begin{definition}[Event]
An \textbf{event} is a \@{subset} of a \@{sample-space}. For example, we can say the event $A$ that the roll of a die is odd is $A = \{1,3,5\},$ the event $B$ that the roll of a die is even is $B = \{2,4,6\}$, and that the event that the roll of a die is less than or equal to 4 is $C = \{1,2,3,4\}.$
\end{definition}

A lot of the rest of basic \@{probability} is just set theory from there:

\begin{itemize}
\item The \@{event} that the roll of a die is not less than or equal to 4 is the \@[complement]{complement-of-a-set} of $C$, written as $C^c = \{5,6\}$.
\item The event that the roll is both even and less than or equal to 4 is the intersection of $A$ and $C$, written as $A \cap C = \{2,4\}.$
\item The event that the roll is either even or less than or equal to 4 is the union of $A$ and $C$, written as $A \cup C = \{1,2,3,4,6\}.$
\end{itemize}

We can also have \@{continuous sample spaces}. For example, if we were to randomly choose a number between 0 and 1 inclusive, the \@{sample space} would be $S = [0, 1]$, and the \@{event} that the outcome is a number whose decimal digits are only 2s would be $\{0.2, 0.22, 0.222, \dots \}.$

\subsection{Probability}

\begin{definition}[Probability]
\emph{Note: This applies to sample spaces of discrete events and is a bit hand wavy.}

We can assign a probability or weight to each sample point in a \@{sample space} by giving it a value ranging from 0 to 1. \@{Events} that are more likely to occur have a probability closer to 1, and events that are less likely to occur have a probability closer to 0. We give a probability to all sample points in a sample space such that the sum of the probabilities of all sample points in a sample space is 1.

Then, the \textbf{probability} of an event $A$ is the sum of the probabilities of all sample points in $A.$ Therefore,

\[ 0 \leq P(A) \leq 1, \quad P(\emptyset) = 0, \quad P(S) = 1. \]
\end{definition}

\begin{definition}[Odds]
\notation{\odds}{\operatorname{odds}}
The \textbf{odds} of an \@{event} $A$ occurring are defined as

	\[ \odds(A) = \frac{P(A)}{P(A^c)} = \frac{P(A)}{1 - P(A)} \]

that is, the \@{probability} $A$ occurs divided by the probability of the \@{complement} of $A$ (i.e., the probability that $A$ does not occur.)
\end{definition}
\begin{note}
\@{Odds} are often written as a \@{ratio} of positive \@{integers}. So we'd say that the odds of drawing an ace from a standard deck of 52 cards is

\[ \odds(\text{Ace}) = \frac{\frac{4}{52}}{\frac{48}{52}} = \frac{4}{48} = \frac{1}{12} = 1 : 12. \]

\end{note}

\begin{definition}[Log Probability]
	A \textbf{log probability} is just a \@{logarithm} of a \@{probability}.
\end{definition}
\begin{note}
Using \@{log probability} has some advantages:
\begin{itemize}
\item It converts multiplication of probabilities into addition, which are faster to compute.
\item It improves \@{numerical stability} when very small probabilities are involved.
\item It eliminates \@{exponentiation}.
\item \@{Logarithms} are \@{concave} which is useful for \@{optimization}.
\end{itemize}
\end{note}

\subsection{Probability Rules}

\begin{definition}[Independence]\synonyms{"independent"}
	The \@{events} $A$ and $B$ are said to be \textbf{independent} if

\[ P(A \cap B) = P(A)P(B). \]
\end{definition}

\begin{theorem}[Addition and Multiplication Rules]\label{probability-rules}
The \@{probability} of $A$ or $B$ is the probability of $A$ plus the probability of $B$ minus the probability of $A$ and $B$ occurring together:

\[ P(A \cup B) = P(A ~ \text{or} ~ B) = P(A) + P(B) - P(A ~ \text{and} ~ B) \]

If $A$ and $B$ are mutually exclusive \@{events}, $P(A \cup B)= P(A) + P(B).$

We have to subtract the overlap between $A$ and $B$ to avoid double counting.

Similarly

\[ P(A \cap B) = P(A) + P(B) - P(A \cup B). \]

The probability of $A$ and $B$ is the probability of $A$ times the probability of $B$ given $A$, or equivalently, the probability of $B$ times the probability of $A$ given $B$.

\[ P(A \cap B) = P(A ~ \text{and} ~ B) = P(B ~ \text{and} ~ A) = P(A) \cdot P(B|A) = P(B) \cdot P(A|B) \]

If $A$ and $B$ are \@{independent}, this reduces to $P(A) \cdot P(B)$.
\end{theorem}

\begin{definition}[Conditional Probability]
From this, we can give the formula for conditional probability. The probability of $A$ given $B$ is

\[ P(A|B) = \frac{P(A ~\text{and} ~ B)}{P(B)}. \]

That is to say, the probability of $A$ occurring given $B$ has occurred is the portion of times $B$ occurs that $A$ also occurs.

If $P(A) = P(A|B)$, then $A$ and $B$ are \@{independent} events, and $P(A) \cdot P(B) = P(A ~ \text{and} ~ B).$
\end{definition}

\end{document}
